\chapter{Revisit 2}

\begin{center}
    {\Huge \bfseries Andi - Japanese Handbook}\\[0.5em]
    {\LARGE \color{boxframe} \textbf{Bab Spesial 2: Kilas Balik Tata Bahasa (Lanjutan)}}
\end{center}
\vspace{1em}

% ================= TOPIC 1 =================
\begin{lessonbox}{1. Frequency words (Seberapa Sering?)}
    Halo lagi, petualang cilik! Kalau mau cerita seberapa sering kamu pergi ke konbini, gunakan kata frekuensi. \textbf{よく (Yoku)} = Sering, \textbf{ときどき (Tokidoki)} = Kadang-kadang.
\end{lessonbox}
\begin{convobox}{1}
    \noindent
    \jpword{わたし}{私}{Watashi} \jpkana{は}{wa} 
    \jpkana{よく}{yoku} 
    \jpkana{コンビニ}{konbini} \jpkana{に}{ni} 
    \jpword{い}{行}{i}\jpkana{きます。}{kimasu.}

    \vspace{1.5em}
    \noindent{\Large \color{articolor} \textbf{Arti:} Saya sering pergi ke minimarket.}
\end{convobox}

% ================= TOPIC 2 =================
\begin{lessonbox}{2. Wanting to do something (Ingin Melakukan Sesuatu)}
    Bunda, kalau si kecil punya keinginan, ajarkan mereka mengganti akhiran \textit{〜masu} menjadi \textbf{〜たい (〜tai)}. \textit{Tabemasu} (makan) menjadi \textit{Tabetai} (ingin makan)!
\end{lessonbox}
\begin{convobox}{2}
    \noindent
    \jpword{すし}{寿司}{Sushi} \jpkana{が}{ga} 
    \jpword{た}{食}{ta}\jpkana{べたいです。}{betai desu.}

    \vspace{1.5em}
    \noindent{\Large \color{articolor} \textbf{Arti:} Saya ingin makan sushi.}
\end{convobox}

% ================= TOPIC 3 =================
\begin{lessonbox}{3. で: Time, Method, Material (Partikel Serbaguna "De")}
    Partikel \textbf{で (De)} sangat sakti! Bisa berarti "dengan alat" (dengan bus), "dari bahan" (dari kayu), atau "batas waktu" (dalam 1 jam).
\end{lessonbox}

\begin{convobox}{3}
    \noindent
    \jpkana{バス}{Basu} \jpkana{で}{de} 
    \jpword{えき}{駅}{eki} \jpkana{に}{ni} 
    \jpword{い}{行}{i}\jpkana{きます。}{kimasu.}

    \vspace{1.5em}
    \noindent{\Large \color{articolor} \textbf{Arti:} Saya pergi ke stasiun DENGAN (menggunakan) bus.}
\end{convobox}

% ================= TOPIC 4 =================
\begin{lessonbox}{4. Casual speech (Bahasa Kasual)}
    Saat ngobrol santai dengan teman baru, buang kata \textit{Desu} dan \textit{Masu}! Ini bikin kamu terdengar ramah dan akrab.
\end{lessonbox}
\begin{convobox}{4}
    \noindent
    \jpkana{これ、}{Kore,} \jpkana{とても}{totemo} \jpkana{おいしい！}{oishii!}

    \vspace{1.5em}
    \noindent{\Large \color{articolor} \textbf{Arti:} Ini, enak banget!}
\end{convobox}

% ================= TOPIC 5 =================
\begin{lessonbox}{5. Casual Speech \#2 (Bertanya Kasual)}
    Untuk bertanya dalam bahasa santai, jangan pakai \textit{〜ka?}. Cukup naikkan nada suaramu di akhir kata seperti orang kaget!
\end{lessonbox}
\begin{convobox}{5}
    \noindent
    \jpword{あした}{明日}{Ashita}\jpkana{、}{,} 
    \jpword{い}{行}{i}\jpkana{く？}{ku?}

    \vspace{1.5em}
    \noindent{\Large \color{articolor} \textbf{Arti:} Besok, kamu pergi (berangkat)?}
\end{convobox}

% ================= TOPIC 6 =================
\begin{lessonbox}{6. Time Expressions (Hari \& Tanggal!)}
    \textbf{Info Penting Turis!} Di Jepang, cara membaca tanggal 1 sampai 10 sangat spesial:
    1 (Tsuitachi), 2 (Futsuka), 3 (Mikka), 4 (Yokka), 5 (Itsuka). 
    Hari ini = \textbf{Kyou}, Besok = \textbf{Ashita}, Kemarin = \textbf{Kinou}.
\end{lessonbox}
\begin{convobox}{6}
    \noindent
    \jpword{きょう}{今日}{Kyou} \jpkana{は}{wa} 
    \jpword{ついたち}{一日}{tsuitachi} \jpkana{です。}{desu.}

    \vspace{1.5em}
    \noindent{\Large \color{articolor} \textbf{Arti:} Hari ini adalah tanggal 1.}
\end{convobox}

% ================= TOPIC 7 =================
\begin{lessonbox}{7. Making Requests (Meminta Tolong)}
    Saat mau difotokan atau minta diambilkan barang, gunakan bentuk \textit{Te} lalu tambah \textbf{ください (Kudasai)}. Artinya "Tolong lakukan...".
\end{lessonbox}
\begin{convobox}{7}
    \noindent
    \jpkana{これ}{Kore} \jpkana{を}{o} 
    \jpword{み}{見}{mi}\jpkana{せて}{sete} 
    \jpword{くだ}{下}{kuda}\jpkana{さい。}{sai.}

    \vspace{1.5em}
    \noindent{\Large \color{articolor} \textbf{Arti:} Tolong perlihatkan (barang) ini.}
\end{convobox}

% ================= TOPIC 8 =================
\begin{lessonbox}{8. Finished/Not Yet (Sudah \& Belum)}
    "Apakah kereta Shinkansen-nya sudah datang?" "Belum!". Gunakan \textbf{もう (Mou)} untuk "Sudah" dan \textbf{まだ (Mada)} untuk "Belum".
\end{lessonbox}
\begin{convobox}{8}
    \noindent
    \jpkana{もう}{Mou} \jpword{た}{食}{ta}\jpkana{べました。}{bemashita.} 
    \jpkana{まだ}{Mada} \jpkana{です。}{desu.}

    \vspace{1.5em}
    \noindent{\Large \color{articolor} \textbf{Arti:} Saya SUDAH makan. / Saya BELUM.}
\end{convobox}

% ================= TOPIC 9 =================
\begin{lessonbox}{9. Change - なる/する (Menjadi...)}
    Jika cuaca berubah dari panas menjadi dingin, pakai kata \textbf{なる (Naru)}. Akhiran \textit{i} berubah jadi \textit{ku}. \textit{Samui} (Dingin) -> \textit{Samuku naru}.
\end{lessonbox}
\begin{convobox}{9}
    \noindent
    \jpword{さむ}{寒}{Samu}\jpkana{くなりました。}{ku narimashita.}

    \vspace{1.5em}
    \noindent{\Large \color{articolor} \textbf{Arti:} (Cuacanya) telah menjadi dingin.}
\end{convobox}

% ================= TOPIC 10 =================
\begin{lessonbox}{10. も: No(thing), even (Tidak Ada Apa-apa!)}
    Kata tanya + \textbf{も (mo)} + Negatif = Kosong/Tidak ada! \\
    Nani (Apa) + Mo + Arimasen = Nani mo arimasen (Tidak ada apa-apa!).
\end{lessonbox}
\begin{convobox}{10}
    \noindent
    \jpword{なに}{何}{Nani} \jpkana{も}{mo} 
    \jpkana{ありません。}{arimasen.}

    \vspace{1.5em}
    \noindent{\Large \color{articolor} \textbf{Arti:} Tidak ada apa-apa.}
\end{convobox}

% ================= TOPIC 11 =================
\begin{lessonbox}{11. Particles: か and を (Atau \& Objek)}
    Kamu ingin jus apel ATAU susu? Gunakan partikel \textbf{か (ka)} di antara pilihanmu. Partikel \textbf{を (o)} menandakan barang yang kamu pilih.
\end{lessonbox}
\begin{convobox}{11}
    \noindent
    \jpkana{ジュース}{Juusu} \jpkana{か}{ka} 
    \jpkana{ミルク}{miruku} \jpkana{を}{o} 
    \jpword{の}{飲}{no}\jpkana{みます。}{mimasu.}

    \vspace{1.5em}
    \noindent{\Large \color{articolor} \textbf{Arti:} Saya minum jus ATAU susu.}
\end{convobox}

% ================= TOPIC 12 =================
\begin{lessonbox}{12. Particles: と and が (Dan \& Subjek Pribadi)}
    Kalau perginya bersama-sama, sambung dengan \textbf{と (to)} yang berarti "Dan". \textbf{が (Ga)} adalah penunjuk siapa yang melakukan kegiatan itu.
\end{lessonbox}
\begin{convobox}{12}
    \noindent
    \jpword{ちち}{父}{Chichi} \jpkana{と}{to} 
    \jpword{はは}{母}{haha} \jpkana{が}{ga} 
    \jpword{い}{行}{i}\jpkana{きます。}{kimasu.}

    \vspace{1.5em}
    \noindent{\Large \color{articolor} \textbf{Arti:} Ayah DAN ibu akan pergi.}
\end{convobox}

% ================= TOPIC 13 =================
\begin{lessonbox}{13. Particles ね、よ、わ、なあ (Partikel Ekspresi)}
    Di akhir kalimat, \textbf{ね (ne)} berarti "ya kan?" (meminta setuju), dan \textbf{よ (yo)} berarti "lho!" (kasih info baru). Bikin bahasamu makin keren!
\end{lessonbox}
\begin{convobox}{13}
    \noindent
    \jpkana{これ、}{Kore,} \jpword{たか}{高}{taka}\jpkana{いですね！}{i desu ne!}

    \vspace{1.5em}
    \noindent{\Large \color{articolor} \textbf{Arti:} Ini, mahal ya!}
\end{convobox}

% ================= TOPIC 14 =================
\begin{lessonbox}{14. Comparisons (Arah \& Membandingkan)}
    \textbf{Info Penting Turis!} Hafalkan arah ini: \textbf{Migi} (Kanan), \textbf{Hidari} (Kiri), \textbf{Massugu} (Lurus). Untuk membandingkan, gunakan \textbf{〜のほうがいい (〜no hou ga ii)}.
\end{lessonbox}
\begin{convobox}{14}
    \noindent
    \jpword{みぎ}{右}{Migi} \jpkana{の}{no} 
    \jpword{ほう}{方}{hou} \jpkana{が}{ga} 
    \jpkana{いいです。}{ii desu.}

    \vspace{1.5em}
    \noindent{\Large \color{articolor} \textbf{Arti:} Sebelah KANAN lebih bagus.}
\end{convobox}

% ================= TOPIC 15 =================
\begin{lessonbox}{15. Amounts (Menyebut Jumlah)}
    Taman bermainnya ramai? Gunakan \textbf{たくさん (Takusan - Banyak)} atau \textbf{すこし (Sukoshi - Sedikit)}.
\end{lessonbox}
\begin{convobox}{15}
    \noindent
    \jpword{ひと}{人}{Hito} \jpkana{が}{ga} 
    \jpkana{たくさん}{takusan} 
    \jpkana{います。}{imasu.}

    \vspace{1.5em}
    \noindent{\Large \color{articolor} \textbf{Arti:} Ada BANYAK orang.}
\end{convobox}

% ================= TOPIC 16 =================
\begin{lessonbox}{16. Lists of things (Mendaftar Banyak Barang)}
    Saat belanja suvenir banyak macamnya, gunakan \textbf{や (ya)} untuk bilang "A, B, dan lain-lain". Berbeda dengan "to" yang menyebutkan semua benda.
\end{lessonbox}
\begin{convobox}{16}
    \noindent
    \jpkana{パン}{Pan} \jpkana{や}{ya} 
    \jpkana{ケーキを}{keeki o} 
    \jpword{か}{買}{ka}\jpkana{います。}{kaimasu.}

    \vspace{1.5em}
    \noindent{\Large \color{articolor} \textbf{Arti:} Saya membeli roti, kue (dan lain-lain).}
\end{convobox}

% ================= TOPIC 17 =================
\begin{lessonbox}{17. (Groups of) People (Kumpulan Orang)}
    Untuk membuat jamak (lebih dari satu), tambahkan \textbf{たち (tachi)}. Watashi (Saya) -> Watashi-tachi (Kami). Kodomo (Anak) -> Kodomo-tachi (Anak-anak).
\end{lessonbox}
\begin{convobox}{17}
    \noindent
    \jpword{わたし}{私}{Watashi}\jpkana{たちは}{tachi wa} 
    \jpword{かぞく}{家族}{kazoku} \jpkana{です。}{desu.}

    \vspace{1.5em}
    \noindent{\Large \color{articolor} \textbf{Arti:} Kami adalah keluarga.}
\end{convobox}

% ================= TOPIC 18 =================
\begin{lessonbox}{18. More て usages (Bentuk "Te" untuk Sifat)}
    Bagaimana bilang "Murah dan Enak"? Ubah kata sifat pertama (Yasui) menjadi \textbf{Yasukute}, lalu sambung!
\end{lessonbox}
\begin{convobox}{18}
    \noindent
    \jpword{やす}{安}{Yasu}\jpkana{くて、}{kute,} 
    \jpkana{おいしいです。}{oishii desu.}

    \vspace{1.5em}
    \noindent{\Large \color{articolor} \textbf{Arti:} Murah dan enak.}
\end{convobox}

% ================= TOPIC 19 =================
\begin{lessonbox}{19. Using から (Karena / Mulai Dari)}
    Partikel \textbf{から (kara)} punya 2 arti hebat: "Dari" (waktu/tempat), atau "Karena" (menjelaskan alasan).
\end{lessonbox}
\begin{convobox}{19}
    \noindent
    \jpword{す}{好}{Su}\jpkana{き}{ki} 
    \jpkana{だからです。}{da kara desu.}

    \vspace{1.5em}
    \noindent{\Large \color{articolor} \textbf{Arti:} Karena (saya) suka.}
\end{convobox}

% ================= TOPIC 20 =================
\begin{lessonbox}{20. Linking sentences (Kata Sambung)}
    Agar ceritamu panjang seperti penduduk lokal, pakai kata sambung! \textbf{でも (Demo)} = Tapi. \textbf{それから (Sorekara)} = Kemudian/Lalu.
\end{lessonbox}
\begin{convobox}{20}
    \noindent
    \jpword{たか}{高}{Taka}\jpkana{いです。}{i desu.} 
    \jpkana{でも、}{Demo,} 
    \jpword{か}{買}{ka}\jpkana{います。}{kaimasu.}

    \vspace{1.5em}
    \noindent{\Large \color{articolor} \textbf{Arti:} Mahal. Tapi, saya membelinya.}
\end{convobox}

% ================= TOPIC 21 =================
\begin{lessonbox}{21. Describing nouns (Menjelaskan Kata Benda)}
    Berbeda dengan bahasa Indonesia, di Jepang penjelasannya diletakkan di DEPAN! "Orang yang memakai baju merah" -> Baju merah dipakai + Orang.
\end{lessonbox}
\begin{convobox}{21}
    \noindent
    \jpword{あか}{赤}{Aka}\jpkana{い}{i} 
    \jpword{ふく}{服}{fuku} \jpkana{を}{o} 
    \jpword{き}{着}{ki}\jpkana{た}{ta} 
    \jpword{ひと}{人}{hito}\jpkana{。}{.}

    \vspace{1.5em}
    \noindent{\Large \color{articolor} \textbf{Arti:} Orang yang memakai pakaian merah.}
\end{convobox}

% ================= TOPIC 22 =================
\begin{lessonbox}{22. Describing actions (Bicara Cara Melakukan)}
    "Berjalan dengan CEPAT". Ubah \textit{Hayai} (Cepat) menjadi \textbf{Hayaku} agar bisa menempel ke kata kerja!
\end{lessonbox}
\begin{convobox}{22}
    \noindent
    \jpword{はや}{早}{Haya}\jpkana{く}{ku} 
    \jpword{ある}{歩}{aru}\jpkana{きます。}{kimasu.}

    \vspace{1.5em}
    \noindent{\Large \color{articolor} \textbf{Arti:} Berjalan dengan cepat.}
\end{convobox}

% ================= TOPIC 23 =================
\begin{lessonbox}{23. Using でしょう (Mungkin / Ya Kan?)}
    Kalau kamu menebak cuaca besok atau meminta persetujuan halus, gunakan \textbf{でしょう (Deshou)}. 
\end{lessonbox}
\begin{convobox}{23}
    \noindent
    \jpword{あした}{明日}{Ashita} \jpkana{は}{wa} 
    \jpword{あめ}{雨}{ame} \jpkana{でしょう。}{deshou.}

    \vspace{1.5em}
    \noindent{\Large \color{articolor} \textbf{Arti:} Besok mungkin akan hujan.}
\end{convobox}

% --- SUMMARY BOX CLOSING ---
\vspace{2em}
\begin{tcolorbox}[
    colback=purple!5!white,
    colframe=purple!60!black,
    boxrule=3pt,
    arc=5mm,
    title={\huge\bfseries \sffamily ⛩️ Penutup: Siap Jelajahi Jepang! ⛩️},
    fonttitle=\color{white},
    attach boxed title to top center={yshift=-4mm},
    boxed title style={colback=purple!60!black, rounded corners},
    left=5mm, right=5mm, top=7mm, bottom=5mm
]
\Large Sempurna! Buku Saku Super Cepat seri 2 ini sudah mencakup hampir seluruh kebutuhan komunikasi dasar untuk bertahan dan bersenang-senang di Jepang! Ingat selalu rumus menanyakan arah, menghitung tanggal, dan meminta tolong. Bawalah keceriaan dan keberanian kalian dalam mempraktikkan bahasa Jepang bersama buku ini. \textbf{また日本でお会いしましょう！ (Mata Nihon de o-ai shimashou - Sampai jumpa lagi di Jepang!)}
\end{tcolorbox}

\newpage